% LSPV-4 REPORT::TRACECLOSE1::2026-10-02T04:49:00-07:00
% !TEX program = xelatex
\documentclass[11pt]{article}
\usepackage[margin=0.88in]{geometry}
\usepackage{amsmath,amssymb}
\begin{document}
\section*{LSPV-4 hostile audit: polynomial transcript functionals are P-closed}

\textbf{STATUS: OPEN.} P versus NP remains open. This report proves a closure fact for nonblack-box polynomial transcript/provenance extraction. It proves no unrestricted SAT lower bound and supplies no nonrelativizing or nonalgebrizing hinge.

\paragraph{Recovery preflight.}
Current GitHub readback showed no verified LSPV-4 report newer than \texttt{2026-10-01\_2046\_LSPV-4\_WEIGHTPAD1.tex}. Preserved-run retrieval did not recover an exact failed scientific body or exact cycle-specific write error for the latest missing post-WEIGHTPAD1 run. The immediate failure classification is therefore \textbf{UNKNOWN}. No missing science was reconstructed or retried.

\paragraph{CORPUS REVIEWED.}
The newest verified LSPV-3 artifact found in its canonical directory is \texttt{2026-10-02\_0230\_LSPV-3\_FACTOR1.tex}. The corpus survey covered current \texttt{AGENTS.md}, root \texttt{main.tex}, active task/governance requirements represented there, the structural SU/vMS/Hadamard branch, reconstruction branch, HEC/CRS-PPC/CCSPC/RCCN/Search--G branch, and the verified LSPV-4 chain through WEIGHTPAD1. The load-bearing gaps remain exact SU/vMS typing, charged acquisition/reconstruction, arbitrary-algorithm transfer, and a nonrelativizing/nonalgebrizing hinge.

\paragraph{Primary-literature preflight.}
Baker--Gill--Solovay (SIAM J. Comput. 4 (1975), 431--442) construct recursive oracles relative to which $P=NP$ and others relative to which $P\ne NP$. Aaronson--Wigderson's algebrization framework strengthens the warning for arithmetization-based methods. Chen--Hu--Ren (arXiv:2511.14038, 2025) establish additional algebrization barriers for circuit lower bounds. Ren--Williams (arXiv:2607.09963, 2026) show that adaptive query rounds and witness lengths can matter quantitatively in a restricted exponential-time circuit-lower-bound setting; this does not transfer to unrestricted SAT runtime.

\paragraph{PROVED/EXACT: polynomial transcript-functional closure.}
Fix a deterministic machine $A$ that halts on length-$n$ inputs in time $T_A(n)\le n^k$. Let $\operatorname{Tr}_A(x)$ be its complete standard step-by-step configuration transcript in a fixed reasonable machine encoding. There is a polynomial $p_A$ such that
\[
 |\operatorname{Tr}_A(x)|\le p_A(n),
\]
and this transcript can be generated in polynomial time by simulating $A$ while recording configurations.

Let $E$ be any deterministic polynomial-time transducer and define
\[
 G_A(x)=E(x,\operatorname{Tr}_A(x)).
\]
Then
\[
 \boxed{A\in P\Longrightarrow G_A\in FP.}
\]
Run $A$ while recording its transcript, then run $E$. The composed cost is polynomial. Hence under the counterfactual $P=NP$, if $A$ is a polynomial-time exact SAT decider,
\[
 \boxed{P=NP\Longrightarrow E(x,\operatorname{Tr}_A(x))\in FP.}
\]

\paragraph{Guard and consequence.}
This does not assert that every property of a polynomial transcript is easy. An unbounded predicate of a polynomial string may be hard. The exact closure is
\[
 \boxed{\text{polynomial transcript length}+\text{polynomial extractor}
 \Longrightarrow\text{polynomial composed computation}.}
\]
Thus moving from black-box SAT answers to the full internal polynomial transcript does not by itself create a superpolynomial resource. A contradiction still needs an independent theorem that the relevant functional cannot be polynomial under the same counterfactual, or a necessity theorem forcing an already-superpolynomial primitive cost. Importing a superpolynomial extractor merely moves the missing lower bound into the extractor.

\paragraph{Hostile wrappers and barriers.}
DEAD/COPY can add irrelevant polynomial trace material; polynomial DELAY remains polynomial; INTERP gives polynomial-overhead universal simulation in a reasonable model; RECOMP regenerates deterministic coordinates; static reread and memoization do not defeat the closure. The proof relativizes: polynomial oracle simulation plus a polynomial oracle extractor remains polynomial relative to the same oracle. It also supplies no nonalgebrizing hinge. Baker--Gill--Solovay and current algebrization barriers therefore remain live guards.

\paragraph{SU/vMS/Hadamard and reconstruction reconciliation.}
No structural claim is strengthened. Fixed-realization Schur/Hadamard, exterior-square, common-kernel, shortening and rank identities remain exact only in their declared models. Exact reconstruction-compatible SU09/vMS invocation, legal Boolean realization, stacked innovation/common-kernel growth, sparse-$Q$ normalization, ACQ1 compilation charges, and SAME\_REALIZATION acquisition remain OPEN or CONDITIONAL. FACTOR1's fixed-answer YES/NO factorization remains compatible with this audit.

\paragraph{HEC/CRS-PPC/CCSPC/RCCN/Search--G reconciliation.}
Restricted no-reread/no-reacquisition state theorems remain exact in their declared models. REREAD1, PCLOSE1-R1, RANDCLOSE1, REACQ1, PADINV1, ISOCAN1 and WEIGHTPAD1 continue to block generic promotion to unrestricted static SAT hardness. TRACECLOSE1 adds that full access to the ordinary polynomial transcript of a hypothetical polynomial solver does not evade polynomial closure when the audit itself is polynomial.

\paragraph{DUPLICATION / CONTRADICTION AUDIT.}
TRACECLOSE1 concerns computation transcripts, not witness multiplicities, so it is not a restatement of WEIGHTPAD1. It is consistent with REACQ1 and PCLOSE1-R1 and makes their composition principle explicit for full polynomial traces. No reviewed corpus theorem supplies the missing independent superpolynomial lower bound for a transcript functional.

\paragraph{PQ / cross-domain classification.}
Schur products, quadratic kernels, explicit linear algebra, and the transcript-composition lemma are \textbf{EXACT OBJECT/THEOREM TRANSFERS}. LWE-style external extraction, collision/special-soundness extraction, provenance semantics, syzygies and coding-theoretic innovation remain \textbf{ANALOGY/CANDIDATE}. Cryptographic hardness assumptions, random oracles, VLSI locality, distributed cuts, online no-reread bounds and proof-system-specific lower bounds remain \textbf{FAILED TRANSFERS} to unrestricted static SAT runtime absent a model-preserving theorem.

\paragraph{WHAT CHANGED RELATIVE TO ROOT MAIN.}
Root main already states that semantic/representation complexity is not forced machine runtime and arbitrary-algorithm acquisition is OPEN. TRACECLOSE1 sharpens this: a polynomial-time functional of the complete ordinary transcript of a polynomial solver is itself polynomial-time computable.

\paragraph{ADOPTED / REJECTED / RETIRED PRIOR COMPONENTS.}
Adopt REREAD1/PCLOSE1-R1/RANDCLOSE1/REACQ1 and exact fixed-code structural identities in their declared scopes. Retain RCCN and causal state theorems only in restricted models. Retire as a generic mechanism
\[
 \boxed{\text{full polynomial transcript}+\text{polynomial audit}
 \Rightarrow\text{superpolynomial SAT resource}.}
\]

\paragraph{STALE OR OVERSTATED CLAIMS FOUND.}
Any statement treating nonblack-box access to the whole polynomial trace as sufficient by itself to escape P-closure is overstated. Such access may be necessary for a future hinge but is not sufficient.

\paragraph{Strict status ledger.}
Transcript generation and polynomial transcript-functional composition are \textbf{PROVED/EXACT}. Baker--Gill--Solovay and cited barrier results are \textbf{KNOWN THEOREM}. ACQ1 and specialized reconstruction ledgers remain \textbf{CONDITIONAL}. A source-specific implementation-invariant necessity theorem forcing a non-P-closed resource remains \textbf{CANDIDATE}. Generic polynomial full-trace audit as a route to superpolynomial SAT cost is \textbf{FALSIFIED/RETIRED}. Exact SU/vMS typing, stacked structural growth, charged acquisition, arbitrary-algorithm transfer/TC4, and the nonrelativizing/nonalgebrizing hinge remain \textbf{OPEN}. Therefore P versus NP remains open.

\paragraph{Next hostile target.}
Write the strongest surviving SAME\_REALIZATION resource as an exact functional $C(A,x)$ and ask whether $C$ is polynomial-time computable from a polynomial execution transcript. If yes, TRACECLOSE1 closes that route under the counterfactual. If no, isolate the extra computational power used by $C$ and prove why every SAT decider must pay for it, then attack that necessity under INTERP, RECOMP, static reread, semantic-oracle replacement and algebraic-oracle replacement.

\end{document}
% END_OF_VERIFIED_REPORT::LSPV-4::TRACECLOSE1
